%%%%%ENTREES%%%%
\begin{tabarg}
\multicolumn{3}{l}{\titretab {Entr\'ees}}\\
\hline
$X_{in}$ & Vecteur, \texttt{float} & Variable d'origine\\
$time_{in}$ & Vecteur, \texttt{double} & Base de temps de $X_{in}$ (secondes)\\
$ecart\_avant$ & Coeff., \texttt{double} & �cart maxi avant (secondes)\\
$ecart\_apres$ & Coeff., \texttt{double} & �cart maxi apr�s (secondes)\\
$toler\_mnq$ & Coeff., \texttt{Int\_32} & Seuil de calcul de la moyenne (nb. �chant). A chaque instant de la base de temps d'arriv�e, 
la moyenne est calcul�e si le nombre d'�chantillons dans l'intervalle de temps est strictement sup�rieur � $toler\_mnq$. Sinon, 
le r�sultat est "valeur manquante". Une valeur de 0 provoque le calcul dans tous les cas (m�me s'il n'y a qu'un seul �chantillon).\\
$time_{out}$ & Vecteur, \texttt{double} & Base de temps d'arriv�e (secondes)\\
$size_{out}$ & Coeff., \texttt{Int\_32} & Nombre de valeurs dans la base de temps 
d'arriv�e (dans adc, utiliser \$Arg6.NbVal)\\
\end{tabarg}

\medskip
%%%%%SORTIES%%%%
\begin{tabarg}
\multicolumn{3}{l}{\titretab {Sorties}}\\
\hline
$X_{out}$ & Vecteur, \texttt{float} & $X_{in}$ moyenn� et �chantillonn� sur
la base de $time_{out}$ (m�me unit� que $X_{in})$\\
\end{tabarg}
